\newpage
\sectiontitle{6}{Rows, open signatures, and generated handlers}
\[
\begin{aligned}
 r&=[(\ell_0,D_0),\ldots,(\ell_{n-1},D_{n-1})],&
 \op{names}(r)&=[\ell_0,\ldots,\ell_{n-1}],\\
 E_r&=[\op{tag}(\ell_0),\ldots,\op{tag}(\ell_{n-1})],&
 \op{tuple}([])&=\unit,\\
 \op{tuple}(t::t_s)&=(t,\op{tuple}(t_s)),&
 d_r&=\op{tuple}([D_0,\ldots,D_{n-1}]),\\
 p_r&:=\lambda c.\,\op{switch}_1\,E_r\,(\lambda\_.\,\op{IDesc}\,I)\,d_r\,c,&
 \rows{r}&:=\code{\sigma}\,E_r\,p_r.
\end{aligned}
\]
Rows are finite lists of literal identities and core description terms.
\[
\begin{aligned}
 \rowok{I}{r}\quad:\!\Longleftrightarrow\quad&
 (\jd{I}{\Set_0})\ \land\ \op{NoDup}(\op{names}(r))\\
 &{}\land\ \forall(\ell,D)\in r.\ \jd{D}{\op{IDesc}\,I}.
\end{aligned}
\]
\[
\ruleof{view_rows}{\rowok{I}{r}\\\jd{D}{\op{IDesc}\,I}\\D\equiv\rows{r}}
 {\view{D}{r}}
\]
\textbf{Open signature assembly.}
\[
 \op{signature}_I(r):=\lambda i.\,\rows{r_i},\qquad
 r_i\text{ is scoped in }\Gamma,i:I.
\]
\[
\ruleof{signature_typing}{\jd{I}{\Set_0}\\
 \op{RowOK}_{\Gamma,i:I}(I,r_i)}
 {\jd{\lambda i.\,\rows{r_i}}{\Defs\,I}}
\]
This derived statement is unproved. Surface $\{R::\op{branches}\,R\}$
uses positive slots: $R\,j$ becomes $\code{\op{var}}\,j$.
The relation starts with those translated schemes, as in the original Rocq presentation.

\textbf{Dispatcher and branch handlers (core terms).}
\[
\begin{aligned}
 M_{I,r,X,Y}&:=\lambda c.\,\interp{p_r\,c}{X}\to Y,\\
 \op{dispatch}^{k}_{I,r,X,Y}(h_s)
 &:=\lambda z.\,
   \bigl(\op{switch}_k\,E_r\,M_{I,r,X,Y}\,
       (\op{tuple}(h_s))\,(\op{fst}\,z)\bigr)(\op{snd}\,z),\\
 \op{retag}_m&:=\lambda x_s.\,(\pos{m},x_s),\\
 \op{exfalso}^{k}_{Y}(d)&:=\lambda x_s.\,\abort_k\,Y\,(d\,x_s).
\end{aligned}
\]
$X$ is the recursive family; $Y$ is the handler result type.
The switch motive depends on the source tag.

\textbf{Row handler relation.}
$\rh{X}{Y}{r'}{r}{h_s}$ aligns source $r$ with target $r'$.
\begin{mathpar}
\ruleof{rh_nil}{}{\rh{X}{Y}{r'}{[]}{[]}}
\\
\ruleof{rh_live}{r'[m]=(\ell,D')\\
 \interp{D}{X}\equiv\interp{D'}{X}\\
 \rh{X}{Y}{r'}{r}{h_s}}
 {\rh{X}{Y}{r'}{(\ell,D)::r}{\op{retag}_m::h_s}}
\\
\ruleof{rh_dead}{\dead{D}{X}{d}\\\rh{X}{Y}{r'}{r}{h_s}}
 {\rh{X}{Y}{r'}{(\ell,D)::r}{\op{exfalso}^{0}_{Y}(d)::h_s}}
\end{mathpar}
A live handler preserves the payload and translates the local tag.
A dead handler requires no occurrence of $\ell$ in $r'$.
Row and endpoint formation checks are supplied by \texttt{ds\_rows} below.

\newpage
\sectiontitle{7}{Certified empty descriptions and description coercions}
$\dead{D}{X}{d}$ generates a witness
$d:\interp{D}{X}\to\bot$.
\begin{mathpar}
\ruleof{dead_bot}{\jd{I}{\Set_0}\\\jd{D}{\op{IDesc}\,I}\\
 \jd{X}{I\to\Set_0}\\D\equiv\code{\bot}}
 {\dead{D}{X}{\id}}
\\
\ruleof{dead_nil}{\jd{I}{\Set_0}\\\jd{D}{\op{IDesc}\,I}\\
 \jd{X}{I\to\Set_0}\\D\equiv\code{\sigma}\,[]\,p}
 {\dead{D}{X}{\lambda z.\,\op{fst}\,z}}
\\
\ruleof{dead_left}{\jd{I}{\Set_0}\\\jd{D}{\op{IDesc}\,I}\\
 \jd{X}{I\to\Set_0}\\D\equiv D_1\dprod D_2\\\dead{D_1}{X}{d}}
 {\dead{D}{X}{\lambda z.\,d(\op{fst}\,z)}}
\\
\ruleof{dead_right}{\jd{I}{\Set_0}\\\jd{D}{\op{IDesc}\,I}\\
 \jd{X}{I\to\Set_0}\\D\equiv D_1\dprod D_2\\\dead{D_2}{X}{d}}
 {\dead{D}{X}{\lambda z.\,d(\op{snd}\,z)}}
\\
\ruleof{dead_choice}{\jd{I}{\Set_0}\\\jd{D}{\op{IDesc}\,I}\\
 \jd{X}{I\to\Set_0}\\\view{D}{r}\\\dr{r}{X}{h_s}}
 {\dead{D}{X}{\op{dispatch}^{0}_{I,r,X,\bot}(h_s)}}
\\
\ruleof{dead_provided}{\jd{I}{\Set_0}\\\jd{D}{\op{IDesc}\,I}\\
 \jd{X}{I\to\Set_0}\\\jd{d}{\interp{D}{X}\to\bot}}
 {\dead{D}{X}{d}}
\\
\ruleof{dr_nil}{}{\dr{[]}{X}{[]}}
\and
\ruleof{dr_cons}{\dead{D}{X}{d}\\\dr{r}{X}{d_s}}
 {\dr{(\ell,D)::r}{X}{d::d_s}}
\end{mathpar}
An explicit empty-payload witness may be checked; non-reduction is not
an emptiness test.

\textbf{Description coercions.}
$\dsub{D}{D'}{X}{q}$ generates
$q:\interp{D}{X}\to\interp{D'}{X}$.
\begin{mathpar}
\ruleof{ds_conv}{\jd{I}{\Set_0}\\
 \jd{D}{\op{IDesc}\,I}\\\jd{D'}{\op{IDesc}\,I}\\
 \jd{X}{I\to\Set_0}\\D\equiv D'}
 {\dsub{D}{D'}{X}{\id}}
\\
\ruleof{ds_dead}{\dead{D}{X}{d}\\
 \jd{I}{\Set_0}\\\jd{D'}{\op{IDesc}\,I}\\\jd{X}{I\to\Set_0}}
 {\dsub{D}{D'}{X}{\op{exfalso}^{0}_{\interp{D'}{X}}(d)}}
\\
\ruleof{ds_rows}{\jd{I}{\Set_0}\\
 \jd{D}{\op{IDesc}\,I}\\\jd{D'}{\op{IDesc}\,I}\\\jd{X}{I\to\Set_0}\\
 \view{D}{r}\\\view{D'}{r'}\\
 \rh{X}{\interp{D'}{X}}{r'}{r}{h_s}}
 {\dsub{D}{D'}{X}{\op{dispatch}^{0}_{I,r,X,\interp{D'}{X}}(h_s)}}
\end{mathpar}

\newpage
\sectiontitle{8}{Type coercions}
Outer-constructor restriction is motivated by
\cite[Section~2]{campos-constructor-subsets}; these coercion rules are the revision.
$\subt{A}{B}{c}$ generates a core function $c:A\to B$.
\begin{mathpar}
\ruleof{su_conv}{\jd{A}{\Set_j}\\\jd{B}{\Set_k}\\A\equiv B}
 {\subt{A}{B}{\id}}
\and
\ruleof{su_trans}{\subt{A}{B}{c}\\\subt{B}{C}{d}}
 {\subt{A}{C}{\lambda x.\,d(c\,x)}}
\\
\ruleof{su_bottom}{\jd{A}{\Set_k}}
 {\subt{\bot}{A}{\lambda z.\,\abort_k\,A\,z}}
\\
\ruleof{su_pi}{\jd{\Pi a:A.\,B}{\Set_j}\\
 \jd{\Pi x:A'.\,B'}{\Set_k}\\
 \subt{A'}{A}{c}\\
 \Gamma,x:A'\vdash B[c\,x/a]\sqsubseteq B'\rightsquigarrow d}
 {\subt{(\Pi a:A.\,B)}{(\Pi x:A'.\,B')}{\lambda f\,x.\,d(f(c\,x))}}
\\
\ruleof{su_close}{\jd{I}{\Set_0}\\
 \jd{F}{\Defs\,I}\\\jd{H}{\Defs\,I}\\\jd{G}{\Defs\,I}\\\jd{i}{I}\\
 \dsub{F\,i}{H\,i}{\cl{G}{G}}{q}}
 {\subt{C_{F,G}(i)}{C_{H,G}(i)}
  {\lambda x.\,\roll\bigl(q(\unroll_{F,G,i}\,x)\bigr)}}
\end{mathpar}
Repeated identical formation premises have been combined.
The recursive definition is the same $G$ on both sides of \texttt{su\_close}.

\textbf{Dependent-function instantiation.}
For element type $S$, instantiate \texttt{su\_pi} as follows:
\[
\begin{gathered}
 \jd{S}{\Set_0},\qquad
 \jd{\Pi a:\op{List}\,S.\,B(a)}{\Set_j},\qquad
 \jd{\Pi x:\op{NonEmpty}\,S.\,B'(x)}{\Set_k},\\
 \subt{\op{NonEmpty}\,S}{\op{List}\,S}{\op{toList}},\\
 \Gamma,x:\op{NonEmpty}\,S
   \vdash B(\op{toList}\,x)\sqsubseteq B'(x)\rightsquigarrow d_x,\\
 (\lambda f\,x.\,d_x(f(\op{toList}\,x))):
   (\Pi a:\op{List}\,S.\,B(a))\to
   (\Pi x:\op{NonEmpty}\,S.\,B'(x)).
\end{gathered}
\]
The source codomain is $B(\op{toList}\,x)$, not $B(x)$.

\textbf{Rocq binder instantiation.}
For de Bruijn terms, $\uparrow^m_c t$ raises indices at cutoff $c$ by $m$;
$[u/0]t$ substitutes for index $0$.
\[
\begin{aligned}
 \op{coercedCodomain}(c,B)
   &:=[\,(\uparrow^1_0 c)\,0/0\,](\uparrow^1_1 B),\\
 \op{piCoercion}(c,d)
   &:=\lambda.\lambda.\,
       (\uparrow^1_1 d)\bigl(1\,((\uparrow^2_0 c)\,0)\bigr).
\end{aligned}
\]
The lifted $d$ remains under the new argument binder; the function binder
is inserted below it. For outer-context variable $c=0$ and $B=1\,0$:
\[
 \op{coercedCodomain}(0,1\,0)=1\,(1\,0).
\]

\newpage
\sectiontitle{9}{Declarative source typing and elaboration}
There is one source judgment:
\[
 \elab{e}{A}{t}.
\]
It relates a source expression $e$, a core type $A$, and its elaboration $t$.
The core judgment remains $\jd{t}{A}$. There is no synthesis/checking split:
$A$ is not assigned an input or output mode, and every source premise below
uses the same relation. Annotations constrain types but are not mode switches.
\begin{mathpar}
\ruleof{el_var}{\wf\\x:A\in\Gamma}{\elab{x}{A}{x}}
\and
\ruleof{el_ann}{\jd{A}{\Set_k}\\\elab{e}{A}{t}}
 {\elab{(e:A)}{A}{t}}
\\
\ruleof{el_app}{\jd{\Pi x:A.\,B}{\Set_k}\\
 \elab{f}{\Pi x:A.\,B}{t}\\\elab{a}{A}{u}}
 {\elab{f\,a}{B[u/x]}{t\,u}}
\\
\ruleof{el_lam}{\jd{\Pi x:A.\,B}{\Set_k}\\
 \Gamma,x:A\vdash e:B\rightsquigarrow t}
 {\elab{\lambda x.\,e}{\Pi x:A.\,B}{\lambda x.\,t}}
\\
\ruleof{el_pair}{\jd{\Sigma x:A.\,B}{\Set_k}\\
 \elab{e}{A}{t}\\\elab{f}{B[t/x]}{u}}
 {\elab{(e,f)}{\Sigma x:A.\,B}{(t,u)}}
\\
\ruleof{el_core}{\jd{t}{A}}{\elab{\op{core}(t)}{A}{t}}
\\
\ruleof{el_conv}{\elab{e}{A}{t}\\\jd{B}{\Set_k}\\A\equiv B}
 {\elab{e}{B}{t}}
\\
\ruleof{el_sub}{\elab{e}{A}{t}\\\subt{A}{B}{c}}
 {\elab{e}{B}{c\,t}}
\\
\ruleof{el_signature}{\jd{I}{\Set_0}\\
 \op{RowOK}_{\Gamma,i:I}(I,r_i)}
 {\elab{\op{signature}\ I\ r}{\Defs\,I}{\lambda i.\,\rows{r_i}}}
\\
\ruleof{el_close}{\jd{I}{\Set_0}\\
 \elab{F}{\Defs\,I}{f}\\\elab{G}{\Defs\,I}{g}}
 {\elab{F(G)}{I\to\Set_0}{\cl{f}{g}}}
\end{mathpar}
For unindexed surface datatypes, the unit application is inserted:
$F(G)$ abbreviates $(\cl{f}{g})\,\unit$.
Dependent substitution uses the elaborated core argument $u$ or component
$t$, not the corresponding source expression.

\textbf{Unrestricted conversion and subsumption.}
\texttt{el\_conv} replaces both former conversion rules without inserting
code. \texttt{el\_sub} may follow \emph{any} source derivation and inserts the
chosen coercion; it is not restricted to variables, applications, or annotated
terms. In particular, lambdas, pairs, constructors, and embedded core terms
can occur directly as application operands.

\textbf{Declarative, not an inference algorithm.}
These rules impose no proof-search order, principal-type property, or
uniqueness of elaboration. Erasing modes from an original derivation produces
a derivation here (also rewriting its clause-compilation subderivations).
The converse is not asserted: unrestricted composition admits new derivations.
Coherence requires separate attention; see Section~13.

\newpage
\sectiontitle{10}{Named constructors and exhaustive cases}
\[
\ruleof{el_constructor}{\jd{I}{\Set_0}\\\jd{F}{\Defs\,I}\\
 \jd{G}{\Defs\,I}\\\jd{i}{I}\\\view{F\,i}{r}\\
 r[n]=(\ell,D)\\\elab{e}{\interp{D}{X_G}}{x_s}}
 {\elab{\ell(e)}{C_{F,G}(i)}{\roll(\pos{n},x_s)}}
\]
The type $C_{F,G}(i)$ in the conclusion fixes the row, local tag, and
recursive payload. It is a parameter of a declarative derivation, not a
checking-mode input. A constructor identity alone need not determine a unique
type; distinct derivations may instantiate $F,G,i$ differently.

\textbf{Clause compilation.}
Let $b_s$ be a list of clauses $(\ell,y.e)$, with $y$ bound in $e$.
$\op{names}(b_s)$ is their list of identities.
$\ecases{X}{Q}{b_s}{r}{h_s}$ compiles one handler for each entry of $r$.
\begin{mathpar}
\ruleof{cases_nil}{}{\ecases{X}{Q}{b_s}{[]}{[]}}
\\
\ruleof{cases_live}{(\ell,y.e)\in b_s\\
 \Gamma,y:\interp{D}{X}\vdash e:Q\rightsquigarrow t\\
 \ecases{X}{Q}{b_s}{r}{h_s}}
 {\ecases{X}{Q}{b_s}{(\ell,D)::r}{(\lambda y.\,t)::h_s}}
\\
\ruleof{cases_dead}{\ell\notin\op{names}(b_s)\\\dead{D}{X}{d}\\
 \ecases{X}{Q}{b_s}{r}{h_s}}
 {\ecases{X}{Q}{b_s}{(\ell,D)::r}{\op{exfalso}^{k}_{Q}(d)::h_s}}
\end{mathpar}
$Q$ is scoped in $\Gamma$, so $y$ does not occur free in $Q$.
An omitted clause needs a checked dead witness.

\textbf{Case typing and elaboration.}
\[
\ruleof{el_case}{\jd{I}{\Set_0}\\\jd{F}{\Defs\,I}\\\jd{G}{\Defs\,I}\\
 \jd{i}{I}\\\elab{e}{C_{F,G}(i)}{x}\\\jd{Q}{\Set_k}\\
 \view{F\,i}{r}\\\op{NoDup}(\op{names}(b_s))\\
 \forall(\ell,y.e_\ell)\in b_s.\ \ell\in\op{names}(r)\\
 \ecases{X_G}{Q}{b_s}{r}{h_s}}
 {\elab{\op{case}_{Q}\ e\ \op{of}\ b_s}{Q}
   {\op{closeCase}_k\,F\,G\,i\,(\lambda\_.\,Q)\,
     (\op{dispatch}^{k}_{I,r,X_G,Q}(h_s))\,x}}
\]
The row view supplies distinct constructor identities. The two clause
premises reject duplicate and foreign identities.

\textbf{Single live clause instantiation.}
For $r=[(\op{cons},K_A)]$, $X=X_{L_A}$, and
$b_s=[(\op{cons},y.\op{core}(\op{fst}\,y))]$:
\[
\begin{gathered}
 \interp{K_A}{X_{L_A}}\equiv A\times\op{List}\,A,\qquad
 \Gamma,y:A\times\op{List}\,A\vdash\op{fst}\,y:A,\\
 \ecases{X_{L_A}}{A}{b_s}{[(\op{cons},K_A)]}
       {[\lambda y.\,\op{fst}\,y]},\\
 \op{case}_{A}\ e\ \op{of}\ b_s
 \rightsquigarrow
 \op{closeCase}_0\,U_A\,L_A\,\unit\,(\lambda\_.\,A)\,
   \bigl(\op{dispatch}^{0}_{\one,r,X_{L_A},A}
      ([\lambda y.\,\op{fst}\,y])\bigr)\,x.
\end{gathered}
\]
Here $\elab{e}{\op{NonEmpty}\,A}{x}$;
$L_A,U_A,K_A$ are defined next. The scrutinee and all live clause bodies use
the same source relation, including its conversion and subsumption rules.
